
\documentclass{article}
\usepackage{colortbl}
\usepackage{makecell}
\usepackage{multirow}
\usepackage{supertabular}

\begin{document}

\newcounter{utterance}

\centering \large Interaction Transcript for game `bbh', experiment `disambiguation\_qa', episode 11 with reflexion{-}Qwen3.5{-}9B.
\vspace{24pt}

{ \footnotesize  \setcounter{utterance}{1}
\setlength{\tabcolsep}{0pt}
\begin{supertabular}{c@{$\;$}|p{.15\linewidth}@{}p{.15\linewidth}p{.15\linewidth}p{.15\linewidth}p{.15\linewidth}p{.15\linewidth}}
    \# & \multicolumn{2}{c}{Player} && \multicolumn{2}{c}{Game Master} \\
    \hline

    \theutterance \stepcounter{utterance}  
    & & & \multicolumn{4}{p{0.6\linewidth}}{
        \cellcolor[rgb]{0.9,0.9,0.9}{
            \makecell[{{p{\linewidth}}}]{
                \texttt{\tiny{[P1$\langle$GM]}}
                \texttt{Clarify the meaning of sentences with ambiguous pronouns.} \\
\\ 
\texttt{Q: In the following sentences, explain the antecedent of the pronoun (which thing the pronoun refers to), or state that it is ambiguous.} \\
\texttt{Sentence: The chief told the counselor that they took the day off.} \\
\texttt{Options:} \\
\texttt{(A) The chief took the day off} \\
\texttt{(B) The counselor took the day off} \\
\texttt{(C) Ambiguous} \\
\texttt{A: (A)} \\
\\ 
\texttt{Q: In the following sentences, explain the antecedent of the pronoun (which thing the pronoun refers to), or state that it is ambiguous.} \\
\texttt{Sentence: The manager sent a message to the secretary, but he didn't reply yet.} \\
\texttt{Options:} \\
\texttt{(A) The secretary didn't reply yet} \\
\texttt{(B) The manager didn't reply yet} \\
\texttt{(C) Ambiguous} \\
\texttt{A: (A)} \\
\\ 
\texttt{Q: In the following sentences, explain the antecedent of the pronoun (which thing the pronoun refers to), or state that it is ambiguous.} \\
\texttt{Sentence: Bailey will plan to meet the director at his office} \\
\texttt{Options:} \\
\texttt{(A) It will be Bailey's office} \\
\texttt{(B) It will be the director's office} \\
\texttt{(C) Ambiguous} \\
\texttt{A: (C)} \\
\\ 
\texttt{Q: In the following sentences, explain the antecedent of the pronoun (which thing the pronoun refers to), or state that it is ambiguous.} \\
\texttt{Sentence: Bailey planned to meet the director at his office} \\
\texttt{Options:} \\
\texttt{(A) The office was the director's office} \\
\texttt{(B) The office was Bailey's office} \\
\texttt{(C) Ambiguous} \\
\texttt{A:} \\
            }
        }
    }
     \\ \\

    \theutterance \stepcounter{utterance}  
    & \multicolumn{4}{p{0.6\linewidth}}{
        \cellcolor[rgb]{0.561,0.737,0.561}{
            \makecell[{{p{\linewidth}}}]{
                \texttt{\tiny{[P1$\rangle$GM]}}
                \texttt{(C)} \\
\\ 
\texttt{**Reasoning:**} \\
\texttt{In the sentence "Bailey planned to meet the director at his office," the pronoun "his" could grammatically refer to either "Bailey" or "the director." Without additional context or specific phrasing (such as "Bailey planned to meet the director at *the director's* office" or "Bailey planned to meet the director at *his {[}Bailey's{]}* office"), it is impossible to determine definitively whose office is being discussed. Therefore, the antecedent is ambiguous.} \\
            }
        }
    }
    & & \\ \\

    \theutterance \stepcounter{utterance}  
    & & & \multicolumn{2}{p{0.3\linewidth}}{
        \cellcolor[rgb]{0.9,0.9,0.9}{
            \makecell[{{p{\linewidth}}}]{
                \texttt{\tiny{[GM$|$GM]}}
                \texttt{Target: (C)} \\
            }
        }
    }
    & & \\ \\

    \theutterance \stepcounter{utterance}  
    & & & \multicolumn{2}{p{0.3\linewidth}}{
        \cellcolor[rgb]{0.9,0.9,0.9}{
            \makecell[{{p{\linewidth}}}]{
                \texttt{\tiny{[GM$|$GM]}}
                \texttt{game\_result = LOSE} \\
            }
        }
    }
    & & \\ \\

    \theutterance \stepcounter{utterance}  
    & & \multicolumn{4}{p{0.6\linewidth}}{
        \cellcolor[rgb]{0.95,0.95,0.95}{
            \makecell[{{p{\linewidth}}}]{
                \texttt{\tiny{[P1$\leftrightarrow$P]}}
                \texttt{Post{-}episode Reflexion (not sent to GM):} \\
\texttt{New plan: I incorrectly identified the sentence as ambiguous because I focused solely on the grammatical possibility of "his" referring to either Bailey or the director. However, in standard English usage, the pronoun "his" in the phrase "meet the director at his office" typically defaults to the subject of the main clause (Bailey) unless context suggests otherwise, making (B) the most likely intended answer rather than (C). Next time, I will apply pragmatic principles of default reference, assuming the pronoun refers to the subject (Bailey) unless there is explicit evidence to the contrary, and select option (B).} \\
            }
        }
    }
    & \\ \\

\end{supertabular}
}

\end{document}
